%-------------------------------------------------------------------------------
%	ABSTRACT PAGE
%-------------------------------------------------------------------------------

\addchaptertocentry{Хураангуй} % Хураангуйг гарчигт нэмэх

\begin{center}
{\scshape\Large \univname\par} % Их сургуулийн нэр
{\scshape\large \facname\par}\vspace{0.5cm} % Их сургуулийн нэр
{\huge\textbf{{Хураангуй}} \par}
\bigskip
{\Large{\ttitle} \par} % Тезисийн нэр
\bigskip

{\normalsize \shortname \par} % Зохиогчийн нэр
\addressname
\end{center}

\bigskip

Энэхүү төгсөлтийн ажилд онлайн дуудлага худалдааны бүрэн цогц систем (full-stack system) хөгжүүлсэн бөгөөд үүнд веб апликейшн (React.js + Vite), мобайл апликейшн (React Native + Expo), мөн RESTful API backend сервер (Node.js + Express.js + MongoDB) багтана. Төслийн гол зорилго нь уламжлалт дуудлага худалдааны аргуудын орон зай, цаг хугацааны хязгаарлалтыг арилгаж, олон платформ дээр ажилладаг (multi-platform), бодит цагийн технологи (real-time) бүхий цахим дуудлага худалдааны орчин бүрдүүлэх явдал юм.

Мобайл апликейшн нь React Native болон Expo framework ашиглан хөгжүүлсэн бөгөөд iOS болон Android төхөөрөмжүүд дээр ажиллах чадвартай. Систем нь бодит цагийн харилцаа (Socket.io), олон төрлийн нэвтрэх арга (Google OAuth, Firebase Phone Auth), AI-р дэмжигдсэн категори санал болгох, мөн Монгол зах зээлд тохирсон 66 төрлийн категоритой. Хэрэглэгч бүр дуудлага худалдааны үйл явцыг шууд хянах, оролцох, санал гаргах боломжтой. WebSocket технологийг ашигласнаар мэдээлэл харилцан солилцох хурд сайжирч, хэрэглэгчийн туршлага (UX) өндөр түвшинд хүрсэн.

Хөгжүүлсэн систем нь 3 үндсэн бүрэлдэхүүн хэсэгтэй: (1) Backend API систем - хэрэглэгч, бараа, дуудлага худалдаа, гүйлгээний удирдлага, RESTful API endpoint-үүд, MongoDB өгөгдлийн сан, бодит цагийн Socket.io холболт; (2) Веб аппликейшн - админы удирдлагын хэсэг, дэлгэрэнгүй дансны тохиргоо, бүрэн функционал бүхий дуудлага худалдааны интерфэйс; (3) Мобайл аппликейшн - iOS болон Android дээр ажилладаг хэрэглэгчийн програм. Хэрэглэгч өөрт тохирсон platform дээр системийг ашиглах боломжтой.

Систем нь хуучин эд зүйлсийг дахин ашиглах, хэн нэгэнд хэрэгцээтэй барааг дамжуулан худалдах нөхцөлийг бүрдүүлдэг тул хог хаягдлыг бууруулах, нөөцийг үр ашигтай ашиглах, байгаль орчинд ээлтэй шийдлийг дэмжиж байгаа. Хувь хүн болон жижиг бизнес эрхлэгчид энэхүү системийг ашигласнаар илүү өргөн зах зээлд хүрэх, ил тод, шуурхай худалдаа хийх боломжтой болж, нийгэм, эдийн засаг, байгаль орчны хувьд бодит эерэг нөлөө үзүүлнэ.





